\documentclass{article}

% packages
\usepackage{amsmath, amsthm, thmtools, amsfonts, amssymb, luacode, catchfile, tikzducks, hyperref, ifthen}
\ifcsname c@kobocompile\endcsname
	\usepackage[a5paper, total={1072pt, 1448pt}, margin=10pt, includeheadfoot]{geometry} % set page margins
\else
	\usepackage[a4paper, margin=50pt, includeheadfoot]{geometry}
\fi
\usepackage[shortlabels]{enumitem}
\usepackage[skip=3pt, indent=0pt]{parskip}

% language
\usepackage[bidi=basic, layout=tabular, provide=*]{babel}
\ifcsname c@english\endcsname
	\babelprovide[main, import]{english}
\else
	\babelprovide[main, import]{hebrew}
	\babelprovide{rl}
\fi
%\babelfont{rm}{Libertinus Serif}
\babelfont{rm}[Renderer=Harfbuzz]{Libertinus Serif}
\babelfont{sf}{Libertinus Sans}
\babelfont{tt}{Libertinus Mono}

% style
\AddToHook{cmd/section/before}{\clearpage}	% Add line break before section
\linespread{1.3}
\setcounter{secnumdepth}{0}		% Remove default number tags from sections, this won't do well with theorems
\AtBeginDocument{\setlength{\belowdisplayskip}{3pt}}
\AtBeginDocument{\setlength{\abovedisplayskip}{3pt}}
\graphicspath{ {../images/} }

% operators
\DeclareMathOperator\cis{cis}
\DeclareMathOperator\Sp{Sp}
\DeclareMathOperator\tr{tr}
\DeclareMathOperator\im{Im}
\DeclareMathOperator\re{Re}
\DeclareMathOperator\diag{diag}
\DeclareMathOperator*\lowlim{\underline{lim}}
\DeclareMathOperator*\uplim{\overline{lim}}
\DeclareMathOperator\rng{rng}
\DeclareMathOperator\Sym{Sym}
\DeclareMathOperator\Arg{Arg}
\DeclareMathOperator\Log{Log}
\DeclareMathOperator\dom{dom}
\DeclareMathOperator\supp{Supp}
\DeclareMathOperator\var{Var}
\DeclareMathOperator\cov{Cov}

% commands
%\renewcommand\qedsymbol{\textbf{מש''ל}}
%\renewcommand\qedsymbol{\fbox{\emoji{lizard}}}
\newcommand{\Aa}[0]{\mathcal{A}}
\newcommand{\Bb}[0]{\mathcal{B}}
\newcommand{\CC}[0]{\mathbb{C}}
\newcommand{\Cc}[0]{\mathcal{C}}
\newcommand{\EE}[0]{\mathbb{E}}
\newcommand{\FF}[0]{\mathbb{F}}
\newcommand{\Ff}[0]{\mathcal{F}}
\newcommand{\Ii}[0]{\mathcal{I}}
\newcommand{\Gg}[0]{\mathcal{G}}
\newcommand{\Ll}[0]{\mathcal{L}}
\newcommand{\Mm}[0]{\mathcal{M}}
\newcommand{\NN}[0]{\mathbb{N}}
\newcommand{\Nn}[0]{\mathcal{N}}
\newcommand{\PP}[0]{\mathbb{P}}
\newcommand{\Pp}[0]{\mathcal{P}}
\newcommand{\QQ}[0]{\mathbb{Q}}
\newcommand{\RR}[0]{\mathbb{R}}
\newcommand{\Rr}[0]{\mathcal{R}}
\newcommand{\Ss}[0]{\mathcal{S}}
\newcommand{\TT}[0]{\mathbb{T}}
\newcommand{\Uu}[0]{\mathcal{U}}
\newcommand{\Vv}[0]{\mathcal{V}}
\newcommand{\Ww}[0]{\mathcal{W}}
\newcommand{\ZZ}[0]{\mathbb{Z}}
\newcommand{\acts}[0]{\circlearrowright}
\newcommand{\explain}[2] {
	\begin{flalign*}
		 && \text{#2} && \text{#1}
	\end{flalign*}
}
\newcommand{\maketitleprint}[0]{ \begin{center}
	%\begin{tikzpicture}[scale=3]
	%	\duck[graduate=gray!20!black, tassel=red!70!black]
	%\end{tikzpicture}	
	\includegraphics[width=6cm]{cover}
\end{center}
}

% theorem commands
\newtheoremstyle{c_remark}
	{}	% Space above
	{}	% Space below
	{}% Body font
	{}	% Indent amount
	{\bfseries}	% Theorem head font
	{}	% Punctuation after theorem head
	{.5em}	% Space after theorem head
	{\thmname{#1}\thmnumber{ #2}\thmnote{ \normalfont{\text{(#3)}}}}	% head content
\newtheoremstyle{c_definition}
	{3pt}	% Space above
	{3pt}	% Space below
	{}% Body font
	{}	% Indent amount
	{\bfseries}	% Theorem head font
	{}	% Punctuation after theorem head
	{.5em}	% Space after theorem head
	{\thmname{#1}\thmnumber{ #2}\thmnote{ \normalfont{\text{(#3)}}}}	% head content
\newtheoremstyle{c_plain}
	{3pt}	% Space above
	{3pt}	% Space below
	{\itshape}% Body font
	{}	% Indent amount
	{\bfseries}	% Theorem head font
	{}	% Punctuation after theorem head
	{.5em}	% Space after theorem head
	{\thmname{#1}\thmnumber{ #2}\thmnote{ \text{(#3)}}}	% head content

\ifcsname c@english\endcsname
	\theoremstyle{plain}
	\newtheorem{theorem}{Theorem}[section]
	\newtheorem{lemma}[theorem]{Lemma}
	\newtheorem{proposition}[theorem]{Proposition}
	\newtheorem*{proposition*}{Proposition}
	%\newtheorem{corollary}[theorem]{אין חלופה עברית}

	\theoremstyle{definition}
	\newtheorem{definition}[theorem]{Definition}
	\newtheorem*{definition*}{Definition}
	\newtheorem{example}{Example}[section]
	\newtheorem{exercise}{Exercise}[section]

	\theoremstyle{remark}
	\newtheorem*{remark}{Remark}
	\newtheorem*{solution}{Solution}
	\newtheorem{conclusion}[theorem]{Conclusion}
	\newtheorem{notation}[theorem]{Notation}
\else
	\theoremstyle{c_plain}
	\newtheorem{theorem}{משפט}[section]
	\newtheorem{lemma}[theorem]{למה}
	\newtheorem{proposition}[theorem]{טענה}
	\newtheorem*{proposition*}{טענה}
	%\newtheorem{corollary}[theorem]{אין חלופה עברית}

	\theoremstyle{c_definition}
	\newtheorem{definition}[theorem]{הגדרה}
	\newtheorem*{definition*}{הגדרה}
	\newtheorem{example}{דוגמה}[section]
	\newtheorem{exercise}{תרגיל}[section]

	\theoremstyle{c_remark}
	\newtheorem*{remark}{הערה}
	\newtheorem*{solution}{פתרון}
	\newtheorem{conclusion}[theorem]{מסקנה}
	\newtheorem{notation}[theorem]{סימון}
\fi

% Questions related commands
\newcounter{question}
\setcounter{question}{1}
\newcounter{sub_question}
\setcounter{sub_question}{1}

\ifcsname c@english\endcsname
	\newcommand{\question}[1][0]{
		\ifthenelse{#1 = 0}{}{\setcounter{question}{#1}}
		\section{Question \arabic{question}}
		\addtocounter{question}{1}
		\setcounter{sub_question}{1}
	}

	\newcommand{\subquestion}[1][0]{
		\ifthenelse{#1 = 0}{}{\setcounter{sub_question}{#1}}
		\subsection{Part \alph{sub_question}}
		\addtocounter{sub_question}{1}
	}
\else
	\newcommand{\question}[1][0]{
		\ifthenelse{#1 = 0}{}{\setcounter{question}{#1}}
		\section{שאלה \arabic{question}}
		\addtocounter{question}{1}
		\setcounter{sub_question}{1}
	}

	\newcommand{\subquestion}[1][0]{
		\ifthenelse{#1 = 0}{}{\setcounter{sub_question}{#1}}
		\subsection{סעיף \localecounter{letters.gershayim}{sub_question}}
		\addtocounter{sub_question}{1}
	}
\fi

% import lua and start of document
\directlua{common = require ('../common')}

\GetEnv{AUTHOR}

% headers
\author{\AUTHOR}
\date\today

\title{פתרון מטלה 1 --- תורת המידה, 80517}

\begin{document}
\maketitle
\maketitleprint[blue]

\question{}
תהי קבוצה $X \ne \emptyset$.

\subquestion{}
תהי $\Aa \subseteq \Pp(X)$ כך ש־$X \in \Aa$ ולכל $E_1, E_2 \in \Aa$ מתקיים $E_1 \setminus E_2 \in \Aa$. \\
נראה ש־$\Aa$ היא אלגברה על $X$.
\begin{proof}
	מתקיים $X \in \Aa$ ולכן גם $X \setminus X = \emptyset \in \Aa$. \\
	אם $E \in \Aa$ אז $X \setminus E = E^C \in \Aa$, כלומר יש סגירות למשלים. \\
	לבסוף נרצה להראות סגירות לאיחוד סופי, נניח ש־$E_1, E_2 \in \Aa$, אז מתקיים,
	\[
		E_1 \cap E_2
		= E_1 \setminus E_2^C
	\]
	אבל $E_2^C \in \Aa$ ולכן $E_1 \cap E_2 \in \Aa$.
	בהתאם גם ${(E_1^C \cap E_2^C)}^C = E_1 \cup E_2 \in \Aa$ ולכן יש סגירות לאיחודים סופיים. \\
	מצאנו כי שלוש התכונות של אלגברה חלות על $\Aa$ ובהתאם היא אלגברה.
\end{proof}

\subquestion{}
תהיינה ${\{ \Aa_n \}}_{n = 1}^{\infty} \subseteq \Pp^2(X)$ אלגברות על $X$ כך שמתקיים $\Aa_n \subseteq \Aa_{n + 1}$ לכל $n < \omega$. \\
נוכיח שגם $\Aa = \bigcup \{ \Aa_n \}$ אלגברה על $X$.
\begin{proof}
	בבירור $\emptyset, X \in \Aa$ ולכן מספיק לבדוק סגירות למשלים ולאיחוד סופי. \\
	אם $A \in \Aa$ אז קיים $n < \omega$ מינימלי כך ש־$A \in \Aa_n$, אבל $\Aa_n$ אלגברה ולכן $A^c \in \Aa_n$ וגם $A^C \in \Aa$. \\
	נניח ש־$A, B \in \Aa$ ובאופן דומה נבחר $n < \omega$ מינימלי כך ש־$A, B \in \Aa_n$, אז מתקיים $A \cup B \in \Aa_n$ ונקבל גם סגירות לאיחוד סופי ב־$\Aa$.
\end{proof}

\subquestion{}
נראה כי איחוד שרשרת עולה של $\sigma$־אלגברות לא בהכרח $\sigma$־אלגברה.
\begin{solution}
	נגדיר $X = \omega_1$ וכן $X_n = n$, וכן $\Aa_n = \langle X_n \rangle$, לכל $n < \omega$.
	מהגדרת ה־$\sigma$־אלגברה הנוצרת נקבל ש־$\Aa_n \subseteq \Aa_{n + 1}$ לכל $n$.
	נגדיר $\Aa = \bigcup \{ \Aa_n \} = \langle \omega \rangle$, ונטען שזו היא $\sigma$־אלגברה.
	אבל $X_n \in \Aa$ לכל $n$ ולכן גם $\bigcup_{n < \omega} X_n = \omega \in \Aa$.
	אבל טענה זו נכונה אם ורק אם קיים $n < \omega$ כך ש־$\omega \in \Aa_n$, וזה כמובן לא יתכן.
\end{solution}

\question{}
תהי $\Bb(\RR)$ $\sigma$־אלגברת בורל ותהי $U \subseteq \RR$ פתוחה.

\subquestion{}
נראה ש־$U$ ניתנת להצגה כאיחוד של אוסף קטעים פתוחים זרים בזוגות.
\begin{proof}
	נניח ש־$U = \bigcup_{\alpha \in I} U_{\alpha}$ חלוקה למחלקת קשירות מסילתית של $U$.
	מתקיים $U_{\alpha} \cap U_{\beta} = \emptyset$ לכל $\alpha \ne \beta \in I$ מהגדרת הקשירות.
	נותר להוכיח שקבוצה קשירה מסילתית ב־$\RR$ היא קטע.

	נניח בשלילה ש־$V$ קשירה מסילתית אך לא קטע ולכן קיימת $a \notin V$ כך שקיימים $b < a < c$ כאשר $b, c \in V$.
	אז בהתאמה $V \cap (-\infty, a)$ וכן $V \cap (a, \infty)$ קבוצות פתוחות המהוות חלוקה של $V$, סתירה.
\end{proof}

\subquestion{}
נראה ש־$|I| \le \omega$.
\begin{proof}
	תהי פונקציית בחירה $c : I \to U$ כך ש־$c\ '' I \subseteq \QQ$, אז $|I| > \omega \implies \omega \le |c(U)| > \omega$, וזו כמובן סתירה.
\end{proof}

\subquestion{}
נסיק ש־$\Bb(\RR)$ נוצרת על־ידי אוסף הקטעים הפתוחים ב־$\RR$.
\begin{proof}
	ידוע כבר כי $\BB = \{ B(x, r) \mid x \in \RR, r > 0 \}$ הוא בסיס לטופולוגיה של $\RR$ ולכן מהווה גם בסיס ל־$\Bb(\RR)$. \\
	נוכל לראות זאת בעוד דרך, כל $U \in \tau(\RR)$ היא איחוד בן־מניה של קטעים פתוחים, ולכן אם נבחר את אוסף כל הטעים הפתוחים נקבל שכל $U$ פתוחה היא איחוד של מספר בן־מניה של איברים משם.
\end{proof}

\question{}
תהי $f : X_1 \to X_2$ פונקציה ותהי $\Mm_2$ $\sigma$־אלגברה על $X_2$.\\
נראה כי $\Mm_1 = \{ f^{-1}(A) \mid A \in \Mm_2 \}$ היא $\sigma$־אלגברה על $X_1$.
\begin{proof}
	$\emptyset, X_2 \in \Mm_2$ ולכן גם $f^{-1}(\emptyset) = \emptyset, f^{-1}(X_2) = X_1 \in \Mm_1$. \\
	נניח ש־$E \in \Mm_1$, אז קיימת $F \in \Mm_2$ כך ש־$E = f^{-1}(F)$ אך גם $F^C \in \Mm_2$ ולכן גם $E^C = {(f^{-1}(F))}^C = f^{-1}(F^C) \in \Mm_1$. \\
	נניח ש־${\{ E_n \}}_{n = 1}^{\infty} \subseteq \Mm_1$ ונסמן ${\{ F_n \}}_{n = 1}^{\infty} \subseteq \Mm_2$ כך ש־$f(E_n) = F_n$ לכל $n < \omega$.
	ידוע כי $\bigcup \{ F_n \} \in \Mm_2$ וכן,
	\[
		f^{-1}(\bigcup \{ F_n \})
		= \bigcup \{ f^{-1}(F_n) \mid n < \omega \}
		\in \Mm_1
	\]
	ונסיק כי יש סגירות לאיחוד בן־מניה וש־$\Mm_1$ היא $\sigma$־אלגברה על $X_1$.
\end{proof}

\question{}
יהי מרחב מטרי $(X, d)$ כך שהטופולוגיה $\tau$ המושרית ממנו היא $\sigma$־אלגברה. \\
נראה שזהו מרחב דיסקרטי.
\begin{proof}
	תהי נקודה $x \in X$ ונגדיר את הקבוצות הפתוחות,
	\[
		U_n = B(x, \frac{1}{n})
	\]
	לכל $n > 1$.
	אז מסגירות לחיתוך בן־מניה גם,
	\[
		\bigcap_{n = 1}^{\infty} U_n
		= \{ x \}
		\in \tau
	\]
	כלומר מצאנו ש־$\tau = \Pp(X)$ כרצוי.
\end{proof}

\question{}
תהי $X$ קבוצה ונניח ש־$\Mm$ היא $\sigma$־אלגברה על $X$ כך ש־$|X| \ge \omega$.

\subquestion{}
נראה ש־$\Mm$ מכילה מספר אינסופי של קבוצות זרות.
\begin{proof}
	נניח ש־${\{ X_n \}}_{n = 1}^{\infty} \subseteq X$ סדרת קבוצות מדידות, ונגדיר סדרה חדשה ברקורסיה.
	נגדיר,
	\[
		Y_1 = \begin{cases}
			X_1 & X_1 \ne \emptyset \\
			X & \text{otherwise}
		\end{cases}
		\quad
		Y_{n + 1}
		= \begin{cases}
			Y_n \cap X_{n + 1} & Y_n \cap X_{n + 1} \ne \emptyset \\
			Y_n \cap X_{n + 1}^C & \text{otherwise}
		\end{cases}
	\]
	אז נקבל ש־$Y_n \ne \emptyset$ וכן ש־$\{ Y_n \}$ שרשרת יורדת ביחס ההכלה, ולבסוף נגדיר $Z_n = Y_n \setminus Y_{n + 1}$, אז נובע ש־$Z_n \cap Z_m = \emptyset$ לכל $n \ne m$, כלומר זוהי קבוצה בת־מניה של קבוצות מדידות זרות בזוגות.
\end{proof}

\subquestion{}
נראה ש־$|\Mm| > \omega$.
\begin{proof}
	נגדיר את הפונקציה $f : 2^{\omega} \to \Mm$ על־ידי,
	\[
		f(g) = \bigcup \{ Z_n \mid n < \omega, g(n) = 1 \}
	\]
	ידוע ש־$\{ Z_n \}$ קבוצה של קבוצות מדידות שונות ולכן איחוד בן־מניה שלהן הוא קבוצה מדידה, ונקבל ש־$f$ מעידה על $|\Mm| = 2^{\omega} > \omega$.
\end{proof}

\end{document}
